\section{Related Work}
\label{section_8_1}

\textbf{Self-Speculative Decoding.} To address the limitations of training-based speculative decoding, a number of prior works have explored training-free approaches. Draft\&Verify~\cite{Draft&Verify} and Swift~\cite{swift_specdec} construct the draft model by skipping a subset of layers from the target model. In contrast, MagicDec~\cite{MagicDec} and QuantSpec~\cite{QuantSpec} construct the draft model by simplifying the KV cache of the target model via pruning and reduced-precision quantization, respectively. In some cases, these approaches additionally introduce separate weights for the draft model. Despite the diversity of self-speculative decoding techniques, their performance gains remain limited in low-batch scenarios. Moreover, when these methods introduce separate weights for the draft model, they incur non-negligible memory capacity overhead, which can be particularly problematic for consumer-grade devices.

\textbf{Systems for Speculative Decoding.} Recent research has actively explored hardware and system designs to accelerate speculative decoding. These works typically assume the existence of distinct target and draft models, and propose heterogeneous architectures tailored to their differing workload characteristics. For instance, some studies~\cite{SpecPIM, LP-spec, ADAPTIVE_DRAFT_LENGTH} employ processing-in-memory (PIM) devices to accelerate the draft model, while others, such as DFVG~\cite{Draft_FPGA_DFVG}, utilize FPGA-based acceleration. While such heterogeneous systems are effective for conventional training-based speculative decoding, they are not well-suited for self-speculative decoding, where the draft model is derived as a subset of the target model.

\textbf{Hardware Accelerators for Format Conversion.} Several prior studies~\cite{ecco, oaken, Deca} have investigated augmenting xPU systems with dedicated encoder and decoder units to support data compression and format transformation. For example, Ecco~\cite{ecco} proposes encoder and decoder designs for GPU-oriented cache compression, while Oaken~\cite{oaken} introduces similar mechanisms for KV cache compression. Unlike these prior efforts, which primarily focus on improving the efficiency of lossy compression, Cassandra proposes encoder and decoder architectures specifically designed to support speculative decoding.

